\section{Introduction}

For $s>1$, the Dirichlet beta function is
\[
 \beta(s)=\sum_{k=0}^{\infty}\frac{(-1)^k}{(2k+1)^s}
         =1-\frac1{3^s}+\frac1{5^s}-\cdots.
\]
\begin{theorem}\label{thm:main}
At least three of $\beta(2),\beta(4),\ldots,\beta(364)$ are irrational.
\end{theorem}

\subsection{Comparison with earlier results}

Rivoal and Zudilin~\cite{RZ03} proved that one of
$\beta(2),\beta(4),\ldots,\beta(14)$ is irrational. The endpoint was reduced
to $12$ by Zudilin~\cite{Z19} and to $10$ by Lai and Zhou~\cite{LZ22}.
OpenAI~\cite[Theorem~1.1]{OAI26} subsequently proved that Catalan's
constant $\beta(2)$ is irrational. Krisman~\cite{K26} proved that at
least one of $\beta(4),\beta(6),\ldots,\beta(36)$ is irrational.

Sorokin~\cite[Theorem~S]{S07} proved that, for every sufficiently large
odd $b$ and every even $a>100b\log b$, at least one of
\[
 \beta(b+1),\beta(b+3),\ldots,\beta(b+a-1)
\]
is irrational. Huttner~\cite[Remark~5]{H12} obtains the same conclusion.
Both results have a computable threshold for $b$, but neither gives
its numerical value. To our knowledge, the explicit count of three
irrational values through $364$ in Theorem~\ref{thm:main} is new.

Asymptotic results give bounds valid for all sufficiently large
endpoints. In~\cite{T26}, we improve the previous
logarithmic lower bounds by proving that, for all sufficiently large
even~$s$, at least $0.21\sqrt{s/\log s}$ of
$\beta(2),\beta(4),\ldots,\beta(s)$ are linearly independent over the rationals.
The proofs of asymptotic results can also yield bounds at specified finite
endpoints, so their consequences at $s=364$ must be checked.
Appendix~\ref{sec:asymptotic-comparisons} makes this comparison for the
finite estimates of Rivoal--Zudilin, Nishimoto and Fischler, and for
the bound in~\cite{T26}. The estimates and parameter families examined
there do not yield the three-value count in Theorem~\ref{thm:main}.

\begin{samepage}
\subsection{Canceling two possible exceptions}

Suppose, for a contradiction, that at most two of the listed beta values
are irrational. We construct three sequences of linear forms. A fixed
integer combination of the three forms at each index cancels the
possible irrational terms, producing a single sequence of rational
numbers. Clearing denominators yields a sequence
of integers tending to zero with infinitely many nonzero terms. This is
impossible: a
nonzero integer has absolute value at least~$1$.\par
\end{samepage}

\begin{samepage}
Three sequences allow us to cancel two coefficients. At each positive
even index $n$, the forms have the shape
\begin{equation}\label{eq:preview}
 r_{n,b}=q_b(n)+\sum_{\substack{2\leq i\leq364\\i\text{ even}}}
                  b^i\rho_i(n)\beta(i),\qquad b=1,2,3,
\end{equation}
with rational coefficients. The coefficient column at $\beta(i)$ is a
scalar multiple of the fixed vector $(1,2^i,3^i)^{\mathsf T}$.
A nonzero integer vector perpendicular to two such vectors gives
weights that cancel the same two beta terms for every $n$.
\end{samepage}

\begin{samepage}
For example, the combination
\[
 H_n=15r_{n,1}-6r_{n,2}+r_{n,3}
\]
has coefficient $\rho_i(n)(15-6\cdot2^i+3^i)$ at $\beta(i)$.
This vanishes for $i=2$ and $i=4$.
\end{samepage}

For any chosen pair, fix such weights and write
$H_n=\sum_{b=1}^3w_br_{n,b}$, with constant term
$q_w(n)=\sum_{b=1}^3w_bq_b(n)$. Sections~\ref{sec:arithmetic}
and~\ref{sec:decay} give a common denominator for the coefficients and
show that multiplying by it preserves decay. The remaining beta values
are rational under our assumption. Multiplying by their fixed common
denominator then gives the required integers tending to zero.

We must still show that $H_n$ is nonzero infinitely often. The constant
term could cancel the remaining beta terms, so we distinguish them by
reduction modulo a large prime $p$. Section~\ref{sec:dilation} compares
the forms at $n=2p$ with those at $n=2$. After multiplication by a
suitable power of $p$, every remaining beta term vanishes modulo $p$,
while the constant term is congruent to a nonzero multiple of $q_w(2)$.
Thus $q_w(2)\ne0$ implies $H_{2p}\ne0$ for all sufficiently large $p$.
Section~\ref{sec:certificate} verifies $q_w(2)\ne0$ for every chosen
pair by an exact determinant calculation. Section~\ref{sec:conclusion}
combines these facts to finish the contradiction.

The construction builds on the rational-function and elimination methods
of Kumar and Tosi~\cite{KT26} and Lai and Yu~\cite{LY20}.
Lai and Sprang~\cite{LS25} also use prime dilation to prove nonvanishing.
Our proof uses partial fractions and elementary prime valuations,
together with Stirling's formula, the prime number theorem in its
least-common-multiple form, and Wilson's and Fermat's congruences.
